\section{Conclusion}
\label{sec:conclusion}

We presented Anomaly-Targeted Retrieval (ATR), a training-free framework for content-relevant intra-video retrieval. Scout scans the complete video, and Sniper then jointly receives the same complete video and a retrieved local clip; ATR is therefore global--local reasoning rather than local-only truncation. On UCF-Crime, ATR enables a 4B model to reach 85.17\% accuracy and an 8B model to reach 85.52\%, outperforming a 32B baseline (83.45\%) at $\sim$10\% of its GPU-hour cost. The gains arise mainly from fewer missed anomalies and must be interpreted with the accompanying FPR trade-off. Cross-dataset, audio-visual, and LLaVA-NeXT experiments show transfer beyond the primary Qwen setup, while prompt and backbone sensitivity define clear applicability boundaries. GT-ATR quantifies substantial headroom from better candidate evidence, whereas the offset diagnostic confirms that Scout should not be described as a calibrated frame-level localizer. ATR is therefore a compute-efficient operating strategy for selected lightweight VLMs, not a universal model enhancer.

% We have presented \textbf{Anomaly-Targeted Retrieval (ATR)}, a training-free framework that reconciles the fundamental tension between local sensitivity and global context in long-form video understanding. By formulating anomaly detection as an \textbf{Intra-Video Retrieval} problem, ATR introduces \textbf{Hierarchical Visual Repetition}---a Video-domain Visual RAG paradigm---to achieve efficient heterogeneous dual-view reasoning. This design preserves bidirectional attention benefits while reducing computational complexity from $O((2L)^2)$ to $O((L+\delta)^2)$.

% On the UCF-Crime standard test set ($N$=290), ATR effectively empowers lightweight models: the 4B model reduces FNR from 26.43\% to 11.43\% (relative reduction 56.8\%), achieving 85.17\% video-level accuracy. The 8B model with ATR reaches 85.52\%, surpassing the 32B Baseline (83.45\%) at approximately 8\% of the GPU$\cdot$h cost. ATR improves recall in 10 out of 13 anomaly categories, with the most substantial gains in concealed anomaly scenarios (Shoplifting +33\%). Large-scale robustness evaluation on the full 1,900-video pool confirms these findings.

% Crucially, ATR's efficacy extends beyond single-domain boundaries. Cross-dataset evaluation on XD-Violence ($N$=800) demonstrates consistent generalization. Furthermore, experiments with the multimodal Qwen2.5-Omni-7B show that ATR transfers seamlessly to audio-visual joint reasoning, achieving a massive F1 improvement from 53.50\% to 80.63\% (+27.13) and substantially breaking through the vision-only modality gap ceiling.

% Furthermore, we identify an \textbf{Over-Alignment Phenomenon} where the 32B model exhibits negative gains under ATR, revealing a critical trade-off between safety alignment and surveillance utility. Our findings provide a critical counter-evidence to the prevailing "scale-is-all-you-need" doctrine in safety-critical domains, demonstrating that lightweight models with targeted retrieval offer a superior efficiency-safety trade-off. By democratizing high-performance video understanding on consumer-grade hardware, ATR paves the way for scalable, privacy-preserving edge surveillance systems.
